%% ============================================================================
\section{Ablation Study}
\label{ch:ablation}
\addcontentsline{lot}{section}{Ablation Study}
%% ============================================================================

Section~\ref{ch:results} evaluates \pads{} as a whole.
This chapter isolates the contribution of its individual mechanisms and probes the
conditions under which the design's assumptions fail, across all three
architectures. Sections~\ref{sec:ablation-components}--\ref{sec:ablation-thresholds}
report three sweeps run to completion this session (39 jobs: 12 component
configurations, 15 scratch-capacity configurations, 12 threshold configurations).
Sections~\ref{sec:ablation-input} and~\ref{sec:ablation-shuffle} were part of the
original evaluation plan but were not run in this session; they are retained as
scoped, un-executed future work rather than populated with invented numbers.
As in Section~\ref{ch:results}, every figure below is a single run ($n=1$);
see that chapter's opening note and Section~\ref{ch:futurework} for why, and
what it means for how these comparisons should be read.

\subsection{Component ablation}
\label{sec:ablation-components}

\pads{} combines three mechanisms: sharding, staging, and adaptive policy
selection. Table~\ref{tab:ablation-components} enables each in turn. The
shard+stage combination (row 4 of the original evaluation plan) is not
expressible in the current implementation -- \code{--data\_policy} is a
single-select flag among \{loose, shard, stage\} -- and is reported here as a
concrete implementation limitation rather than a missing run.

\begin{table}[!ht]
\centering
\begin{tabular}{llccc}
\toprule
Architecture & Configuration & GPU util.\ (\%) & IoU (pooled) \\
\midrule
DeepLabV3+ & None (loose)            & 61.0 & 0.3707 \\
           & Shard only              & \textbf{84.1} & 0.3569 \\
           & Stage only              & 67.6 & 0.3582 \\
           & Full \pads{} (adaptive) & 79.2 & 0.3459 \\
\midrule
MA-Net     & None (loose)            & 84.6 & 0.3695 \\
           & Shard only              & 86.8 & 0.3527 \\
           & Stage only              & \textbf{87.0} & 0.3602 \\
           & Full \pads{} (adaptive) & 83.7 & 0.3536 \\
\midrule
SegFormer  & None (loose)            & 42.1 & 0.3647 \\
           & Shard only              & 50.7 & 0.3735 \\
           & Stage only              & \textbf{54.3} & 0.3689 \\
           & Full \pads{} (adaptive) & 41.0 & 0.3757 \\
\bottomrule
\end{tabular}
\caption{Component ablation, GPU utilisation and IoU by policy component, all
three architectures. Shard+stage combined is not expressible in the current
implementation (see text).}
\label{tab:ablation-components}
\end{table}

\begin{figure}[!ht]
\centering
\includegraphics[width=0.85\linewidth]{ablation-components.png}
\caption{GPU utilisation by policy component, all three architectures.}
\label{fig:ablation-components}
\end{figure}

Shard alone beats \code{loose} on GPU utilisation for every architecture -- a
clean, consistent result, and for DeepLabV3+ shard even beats stage (84.1\% vs.\
67.6\%). This confirms the underlying staging mechanism has real value when
applied. The ``full \pads{}'' row, however, is not architecturally distinct from
the loose row in most of these runs: as established in
Section~\ref{sec:res-profiler-bias}, the adaptive profiler selected \code{loose}
in every one of these three full-\pads{} runs, so row~4 is effectively row~1
re-measured under the profiler rather than a genuinely different configuration.
The component ablation therefore separates two questions that a single ``does
\pads{} help'' number would conflate: the staging primitive itself helps
(rows~2--3 beat row~1 in GPU utilisation on every architecture), and the
mechanism that is supposed to decide when to invoke it does not currently invoke
it under these conditions.

\subsection{Scratch capacity: the regime the design targets}
\label{sec:ablation-capacity}

The Bing 1k corpus fits comfortably in node-local scratch, which is the least
interesting case for an adaptive stager: a fixed policy of ``copy everything up
front'' suffices. The design is aimed at the regime where it does not fit, and
that regime is reached here by synthetically constraining the scratch budget.

\begin{table}[!ht]
\centering
\begin{tabular}{lccc}
\toprule
Scratch budget & DeepLabV3+ & MA-Net & SegFormer \\
 & GPU util.\ (\%) / IoU & GPU util.\ (\%) / IoU & GPU util.\ (\%) / IoU \\
\midrule
100\% & 17.3 / 0.3383 & 86.9 / 0.3602 & 51.1 / 0.3648 \\
 50\% & 81.6 / 0.3547 & 78.6 / 0.3527 & 56.0 / 0.3716 \\
 25\% & 24.2 / 0.3463 & 85.2 / 0.3602 & 54.0 / 0.3641 \\
 10\% & 87.2 / 0.3564 & 85.1 / 0.3527 & 53.8 / 0.3656 \\
  0\% (staging unavailable) & 79.8 / 0.3646 & 85.8 / 0.3527 & 52.0 / 0.3686 \\
\bottomrule
\end{tabular}
\caption{Behaviour as scratch capacity is constrained. Method~4 of
Section~\ref{sec:evaluation-plan} is unavailable below 100\%.}
\label{tab:ablation-scratch}
\end{table}

\begin{figure}[!ht]
\centering
\includegraphics[width=0.85\linewidth]{ablation-scratch-capacity.png}
\caption{GPU utilisation vs.\ scratch-capacity budget, all three architectures.
DeepLabV3+'s 100\%-capacity run is the worst-performing configuration in the
entire sweep.}
\label{fig:ablation-scratch}
\end{figure}

This is the experiment that was intended to distinguish a contribution from an
observation, and the result is not the graceful-degradation curve the design
predicts. MA-Net and SegFormer are both close to flat across the whole sweep --
consistent with graceful degradation, but also consistent with the scratch
budget simply not being the binding constraint for either architecture at this
corpus size. DeepLabV3+ shows the opposite of the predicted direction: 100\%
scratch capacity produces the \emph{lowest} GPU utilisation in the entire sweep
(17.3\%), well below even the 0\%-capacity (staging-unavailable) row. Staging
overhead is not merely failing to help at full capacity for DeepLabV3+; it is
actively worse than not staging at all. IoU is comparatively stable across the
sweep for every architecture (0.338--0.372), so this is a throughput effect, not
an accuracy one. The honest reading is that ``more caching capacity is better''
does not hold unconditionally in this data, and the mechanism -- whether it is
staging-copy overhead specific to DeepLabV3+'s I/O pattern, or an interaction
with the profiler-bias defect of Section~\ref{sec:res-profiler-bias} that this
sweep did not control for -- is an open question rather than one this experiment
settles.

\subsection{Input size and the CORONA variant}
\label{sec:ablation-input}

Section~\ref{sec:dataset} argued that Bing 1k is the least favourable variant for
a staging argument, since larger tiles increase bytes per sample and therefore
shift the compute-to-I/O ratio. That argument was scoped to be tested against
Bing 2k and Bing+CORONA 1k variants, following Soroush et al.~\cite{Soroush20},
Vadineanu et al.~\cite{Vadineanu24} and Pistola et al.~\cite{Pistola25}, who each
train on CORONA imagery for archaeological site detection in the wider
Mesopotamian region. This ablation was not run in this session -- the
component, scratch-capacity and threshold sweeps above, across three
architectures, were the priority given available cluster time -- and is retained
as scoped future work (Section~\ref{ch:futurework}) rather than populated here.

\subsection{Classifier threshold sensitivity}
\label{sec:ablation-thresholds}

The classification rule of Section~\ref{ch:methodology} depends on a threshold
$\gamma_d$, set by default to 1.5. Values chosen by the designer and never
varied are a standing weakness in any adaptive system, and the sensitivity is
reported here across all three architectures.

\begin{table}[!ht]
\centering
\begin{tabular}{llcccc}
\toprule
Architecture & $\gamma_d$ & Measured ratio & Policy selected & IoU & GPU util.\ (\%) \\
\midrule
DeepLabV3+ & 1.00 & 0.64 & loose & 0.3569 & 64.6 \\
           & 1.25 & 1.23 & loose & 0.3510 & 86.3 \\
           & 1.50 (default) & 0.65 & loose & 0.3513 & 61.6 \\
           & 2.00 & 1.17 & loose & 0.3442 & 83.0 \\
\midrule
MA-Net     & 1.00 & 1.03 & \textbf{shard} & 0.3553 & 86.4 \\
           & 1.25 & 0.69 & loose & 0.3536 & 80.1 \\
           & 1.50 (default) & 1.18 & loose & 0.3720 & 85.1 \\
           & 2.00 & 0.68 & loose & 0.3536 & 87.0 \\
\midrule
SegFormer  & 1.00 & 1.23 & \textbf{shard} & 0.3527 & 49.6 \\
           & 1.25 & 1.22 & loose & 0.3682 & 41.2 \\
           & 1.50 (default) & 1.32 & loose & 0.3738 & 42.5 \\
           & 2.00 & 1.20 & loose & 0.3736 & 43.0 \\
\bottomrule
\end{tabular}
\caption{Sensitivity of the policy decision to the classification threshold, all
three architectures. Only 2 of 12 runs ever left \code{loose}, both at the
single lowest threshold tested and both right at the ratio boundary. These
measured ratios predate the correction of Section~\ref{sec:res-profiler-bias};
under the corrected measurement, every full-window ratio is 10.3--20.2, which
would clear every threshold tested here.}
\label{tab:ablation-thresholds}
\end{table}

\begin{figure}[!ht]
\centering
\includegraphics[width=\linewidth]{ablation-threshold-sensitivity.png}
\caption{Measured ratio vs.\ decision threshold $\gamma_d$, all three
architectures. Circled points switched policy away from \code{loose}.}
\label{fig:ablation-thresholds}
\end{figure}

The relevant question is not which threshold maximises throughput on this
workload -- that would be tuning on the test set -- but how wide the band is
over which the decision is unchanged, and the answer is: very wide, almost the
entire tested range. The policy decision is stable at \code{loose} across three
of four thresholds for every architecture, switching only at the single most
aggressive value tested ($\gamma_d=1.0$), and even then for only two of three
architectures. A decision this insensitive to its own designer-chosen threshold
is not evidence of a well-calibrated decision boundary; combined with the finding
of Section~\ref{sec:res-profiler-bias}, it is evidence that the measured ratio
feeding the decision was itself biased low across this entire sweep, so the
threshold's apparent robustness is partly an artefact of the input it was
comparing against being wrong by roughly an order of magnitude rather than of the
threshold being well-chosen.

\subsection{Shuffling randomness}
\label{sec:ablation-shuffle}

Shard-based reading trades exact random access for sequential throughput. This
ablation was scoped to vary a shuffle-buffer size and establish where, if
anywhere, reduced shuffling randomness begins to affect segmentation quality, in
the manner of the standard WebDataset shuffle-buffer pattern~\cite{Aizman19}.
As documented in Section~\ref{sec:res-accuracy}, \pads{}'s shard reader opens
archive members by name rather than as a sequential stream, so this concern does
not apply to the implementation evaluated in this thesis; the ablation is
retained as future work for completeness against the general class of
shard-based readers this thesis positions itself against, but was not run this
session and is not needed to support the accuracy claims made about \pads{}
specifically.
